\section{An exact colliding control and the scope of the mechanism}\label{sec:control}

The instrument is source independent. Its application to a source uses
positivity, compact-parameter square-root regularity and complete moment
bounds. Consequently the generic identities must be tested on a source
that meets those requirements and nevertheless collides. The following
recorded reflected control \cite[Section 4, (9)--(12)]{ReflectedControl}
provides such a test with an exact event inside the physical slab.

\begin{theorem}[A smooth admissible source with a double event]\label{thm:control}
Let
\begin{align}
 g(v)&=\ee^{-v^4}+\tfrac12\ee^{-(v-4)^4}
                     +\tfrac12\ee^{-(v+4)^4},\label{eq:gcontrol}\\
 \Phi_*(v)&=c\ee^{-v^2/4}g(v),\qquad
 c=\left(16\int_0^\infty\cosh(v)\ee^{-v^2/4}g(v)\dd v\right)^{-1}>0.
 \label{eq:controlsource}
\end{align}
This source admits the same complete angular preparation and count
instrument with its own source normalization. It has all the fixed
square-root and moment domains used above. Its transform has a nondegenerate
double zero at $\lambda_*=1/4$, $k_*=\pi/4$, corresponding to laboratory
time $t_*=1/(4\gamma)$.
\end{theorem}
\begin{proof}
The source is strictly positive, even and smooth. Each derivative of a
summand of $g$ is that summand times a polynomial. Positivity of the mixture
therefore bounds $|g^{(\ell)}/g|$ by a fixed polynomial in $|v|$.
The finite ratio recurrence \eqref{eq:sqrtrecurrence} bounds every derivative
of $\sqrt g$ by a polynomial times $\sqrt g$. For $|v|\ge8$,
$|v-b|\ge|v|/2$ for $b=0,\pm4$, so
\begin{equation}\label{eq:controltail}
 g(v)\le2\ee^{-v^4/16}.
\end{equation}
Thus $\ee^{(\lambda-1/4)v^2/2}\sqrt g$ and its fixed mixed derivatives
have uniformly bounded Schwartz seminorms on each compact real $\lambda$
interval. All polynomial and exponential moments of the source are finite.
The normalization has a positive compact-interval lower bound.
With the same $R$, \eqref{eq:fullmoment} pays the full radial and angular
domains. The instrument proofs are unchanged.

At $\lambda=1/4$, ordinary full-line Fourier translation, followed by
evenness, gives
\begin{equation}\label{eq:controlfactor}
 H^*_{1/4}(k)=cF(k)(1+\cos4k),\qquad
 F(k)=\int_0^\infty\ee^{-v^4}\cos(kv)\dd v.
\end{equation}
Let $M_r=\int_0^\infty v^r\ee^{-v^4}\dd v$.
Integration of $(v\ee^{-v^4})'$ gives $M_4=M_0/4$.
Cauchy--Schwarz implies $M_2^2\le M_0M_4$, hence $M_2\le M_0/2$.
The global inequality $\cos y\ge1-y^2/2$ now gives
\begin{equation}\label{eq:Fpositive}
 F(k)\ge M_0(1-k^2/4)>0\qquad(|k|<2).
\end{equation}
At $k_*=\pi/4$ the factor $1+\cos4k$ and its first derivative vanish,
while its second derivative is $16$. Therefore
\begin{equation}\label{eq:controlcollision}
 H^*_{1/4}(k_*)=(H^*_{1/4})'(k_*)=0,\qquad
 (H^*_{1/4})''(k_*)=16cF(\pi/4)>0.
\end{equation}
Finally $1/2-\gamma t_*=1/4$ gives the stated interior physical time.
Every step is analytical, including the positivity of the curvature.
\end{proof}

The control has a smooth state, positive success mass, all count outputs,
the covariance law, the finite-sector transfer laws and the complete signed
sum in \eqref{eq:conserved}. It satisfies the local collision description
precisely. It fails the canonical integer-square arithmetic expansion and
the canonical theta tail formula; its own super-Gaussian bounds pay the
generic preparation requirements. Hence the theorem excludes an argument
based solely on those generic dynamical properties on a domain containing
this witness. It does not exclude source-sensitive analysis of the unchanged
theta source.

\subsection{What the other controls test}
The wider repository control collection \cite{CollisionControls} separates
premises rather than providing one universal counterexample class:
\begin{center}\small
\begin{tabularx}{\textwidth}{@{}>{\raggedright\arraybackslash}p{0.25\textwidth}Y@{}}
\toprule
Control family & Relevant distinction \\\midrule
Euler-factor sources & Fail the unchanged canonical arithmetic source; applicable positive smooth preparation and Fourier/filtering identities can survive. \\
Weighted-square and finite-prime changes & Complete unit weights or full canonical reflection fail. Finite endpoint matches do not give the smooth full-line source domains. \\
Reflected smooth and positive-primitive sources & Exact theta arithmetic is absent. Positivity, reflection and the applicable generic transport laws need not prevent collisions. \\
Atomic or nonsmooth finite profiles & Fourier moments can exist while the present smooth square-root and graph-domain requirements fail or need a different construction. \\
\bottomrule
\end{tabularx}
\end{center}
Only Theorem~\ref{thm:control}'s preparation and event are checked afresh
here; this table does not claim a new domain audit of every family. A finite
collision-control estimate should be tested where the relevant witness lies,
not only in a different collision-free region.

The analytical gain obtained so far is a complete representation and a
precise collision dictionary. Equations~\eqref{eq:covariance},
\eqref{eq:dj}, \eqref{eq:generate} and \eqref{eq:sectorcurvature} are also
available by differentiating the same source integral. Their reorganization
has not yet supplied an additional signed margin. A new mechanism must
identify a consequence of the canonical source that survives its deformation
and distinguishes the relevant controls quantitatively or structurally.
